\section{ANEXO C. DISTRIBUCIÓN Y PUBLICACIÓN EN EL REGISTRO DE PAQUETES
PYPI}\label{anexo-c.-distribuciuxf3n-y-publicaciuxf3n-en-el-registro-de-paquetes-pypi}

El ciclo de desarrollo de la librería propuesta culmina con su fase de
distribución, permitiendo que las herramientas de análisis de seguridad
implementadas sean accesibles e integrables por la comunidad de
desarrollo y auditoría de \emph{smart contracts}. Para asegurar una
distribución estandarizada y eficiente dentro del ecosistema Python, se
ha seleccionado el índice oficial de paquetes PyPI (\emph{Python Package
Index}). La gestión de este proceso se unifica bajo la herramienta
\texttt{uv}, garantizando la consistencia desde la compilación de los
artefactos hasta su publicación definitiva. \#\# ANEXO C.1.
Configuración de Metadatos del Proyecto (pyproject.toml)

El paso previo indispensable para la distribución consiste en la
definición inequívoca de los metadatos y la especificación del sistema
de construcción (\emph{build system}) en el archivo de configuración
\texttt{pyproject.toml}, ubicado en la raíz del paquete. Este
procedimiento se rige bajo los estándares modernos de empaquetado de
Python (PEP 517 y PEP 621).

Para este proyecto, se ha adoptado \texttt{hatchling} como \emph{build
backend}, debido a su ligereza, velocidad y compatibilidad nativa con
las especificaciones del ecosistema actual. A continuación, se presenta
la estructura de configuración requerida para delimitar las propiedades
de la librería \texttt{evmaudit}:

\begin{Shaded}
\begin{Highlighting}[]
\KeywordTok{[project]}
\DataTypeTok{name} \OperatorTok{=} \StringTok{"evmaudit"}
\DataTypeTok{version} \OperatorTok{=} \StringTok{"0.1.0"}
\DataTypeTok{description} \OperatorTok{=} \StringTok{"Add your description here"}
\DataTypeTok{readme} \OperatorTok{=} \StringTok{"README.md"}
\DataTypeTok{authors} \OperatorTok{=} \OperatorTok{[}
    \OperatorTok{\{ }\DataTypeTok{name}\OperatorTok{ =} \StringTok{"Daniel Rovira Martínez"}\OperatorTok{, }\DataTypeTok{email}\OperatorTok{ =} \StringTok{"pardalaco@gmail.com"}\OperatorTok{ \},}
    \OperatorTok{\{ }\DataTypeTok{name}\OperatorTok{ =} \StringTok{"Paula Suárez Prieto"}\OperatorTok{, }\DataTypeTok{email}\OperatorTok{ =} \StringTok{"Paula Suárez Prieto"}\OperatorTok{ \},}
    \OperatorTok{\{ }\DataTypeTok{name}\OperatorTok{ =} \StringTok{"Adrián Moreno Martín"}\OperatorTok{, }\DataTypeTok{email}\OperatorTok{ =} \StringTok{"adrimore2@gmail.com"}\OperatorTok{ \}}
\OperatorTok{]}
\DataTypeTok{requires{-}python} \OperatorTok{=} \StringTok{"\textgreater{}=3.12"}
\DataTypeTok{dependencies} \OperatorTok{=} \OperatorTok{[}
    \StringTok{"eth\textgreater{}=0.0.1"}\OperatorTok{,}
    \StringTok{"mythril\textgreater{}=0.24.8"}\OperatorTok{,}
    \StringTok{"setuptools\textless{}70.0.0"}\OperatorTok{,}
    \StringTok{"slither{-}analyzer\textgreater{}=0.9.2"}\OperatorTok{,}
    \StringTok{"solc{-}select\textgreater{}=1.2.0"}\OperatorTok{,]}

\KeywordTok{[project.scripts]}
\DataTypeTok{evmaudit} \OperatorTok{=} \StringTok{"evmaudit.main:main"}

\KeywordTok{[build{-}system]}
\DataTypeTok{requires} \OperatorTok{=} \OperatorTok{[}\StringTok{"uv\_build\textgreater{}=0.11.15,\textless{}0.12.0"}\OperatorTok{]}
\DataTypeTok{build{-}backend} \OperatorTok{=} \StringTok{"uv\_build"}
\end{Highlighting}
\end{Shaded}

Los clasificadores (\emph{classifiers}) incluidos permiten categorizar
la librería dentro del índice público, facilitando su indexación en
función de la licencia de código abierto seleccionada, la compatibilidad
del sistema operativo y las versiones soportadas del intérprete Python.

\subsection{ANEXO C.2. Proceso de Compilación del
Paquete}\label{anexo-c.2.-proceso-de-compilaciuxf3n-del-paquete}

Una vez validados los metadatos, se procede a la generación de los
archivos de distribución. Este proceso transforma el código fuente
estructurado en el directorio \texttt{src/} en artefactos instalables e
independientes del entorno de desarrollo.

La ejecución del comando unificado de compilación abstrae la complejidad
de las herramientas tradicionales:

\begin{Shaded}
\begin{Highlighting}[]
\ExtensionTok{uv}\NormalTok{ build}
\end{Highlighting}
\end{Shaded}

Este comando genera de forma nativa dos tipos de distribuciones estándar
dentro del directorio \texttt{dist/}:

\begin{itemize}
\tightlist
\item
  \textbf{Distribución de código fuente (\emph{sdist} o}
  \texttt{.tar.gz}\textbf{):} Un archivo comprimido que contiene el
  código fuente original y los archivos de configuración, actuando como
  respaldo para plataformas o configuraciones no previstas.
\item
  \textbf{Distribución binaria compilada (\emph{wheel} o}
  \texttt{.whl}\textbf{):} El formato de empaquetado moderno que permite
  una instalación directa y optimizada en el sistema destino, evitando
  la necesidad de compilar el código en la máquina del usuario final.
\end{itemize}

En contextos de desarrollo complejos donde el paquete forma parte de un
repositorio principal o arquitectura de \emph{workspace} (como el
entorno de trabajo \texttt{TFM-UNIR}), \texttt{uv} permite gestionar la
compilación de manera localizada. Para forzar la construcción exclusiva
del subproyecto desde su propio directorio y evitar colisiones en la
raíz global, se aplica la opción de empaquetado específico:

\begin{Shaded}
\begin{Highlighting}[]
    \ExtensionTok{uv}\NormalTok{ build }\AttributeTok{{-}{-}package}\NormalTok{ evmaudit}
\end{Highlighting}
\end{Shaded}

\subsection{ANEXO C.3. Seguridad y Autenticación en la
Publicación}\label{anexo-c.3.-seguridad-y-autenticaciuxf3n-en-la-publicaciuxf3n}

La publicación de código en repositorios públicos exige mecanismos
estrictos de control de acceso para prevenir vectores de ataque basados
en la cadena de suministro (\emph{supply chain attacks}). Por razones de
seguridad, se desestima el uso de contraseñas de usuario tradicionales
en favor de la autenticación basada en \textbf{Tokens de API.}

El proceso de despliegue requiere la obtención de un \emph{token} con
prefijo \texttt{pypi-} generado desde el panel de control de PyPI. En la
primera interacción, el alcance del \emph{token} se configura de manera
global; no obstante, una vez realizada la primera subida con éxito, la
buena práctica metodológica dicta restringir los permisos del token de
manera exclusiva al ámbito del paquete \texttt{evmaudit}, minimizando
así la superficie de exposición en caso de compromiso de la credencial.

\subsection{ANEXO C.4. Ejecución del
Despliegue}\label{anexo-c.4.-ejecuciuxf3n-del-despliegue}

Con los artefactos ubicados en el directorio \texttt{dist/} y las
credenciales expedidas, se procede a la transferencia segura hacia los
servidores de PyPI. El comando \texttt{uv\ publish} automatiza la
verificación de integridad mediante \emph{hashes} criptográficos y
realiza la subida en un único paso:

\begin{Shaded}
\begin{Highlighting}[]
\ExtensionTok{uv}\NormalTok{ publish}
\end{Highlighting}
\end{Shaded}

Durante el flujo interactivo en la línea de comandos, el sistema
requiere la introducción del identificador genérico
\texttt{\_\_token\_\_} en el campo de usuario, seguido de la clave
alfanumérica del token de API en el campo de contraseña. Con el objetivo
de optimizar este flujo en entornos de Integración Continua (CI/CD) o
evitar la inserción manual recurrente, es posible exportar temporalmente
la credencial en el entorno de la terminal actual:

\begin{Shaded}
\begin{Highlighting}[]
\BuiltInTok{export} \VariableTok{UV\_PUBLISH\_TOKEN}\OperatorTok{=}\StringTok{"pypi{-}tu{-}token{-}aqui"}
\ExtensionTok{uv}\NormalTok{ publish}
\end{Highlighting}
\end{Shaded}

Tras la finalización exitosa del proceso, el paquete queda registrado
globalmente, permitiendo su incorporación inmediata en otros proyectos
mediante los gestores tradicionales del ecosistema:

\begin{Shaded}
\begin{Highlighting}[]
\ExtensionTok{pip}\NormalTok{ install evmaudit}
\end{Highlighting}
\end{Shaded}

O bien, aprovechando los beneficios de rendimiento de la herramienta
unificada del proyecto:

\begin{Shaded}
\begin{Highlighting}[]
\ExtensionTok{uv}\NormalTok{ add evmaudit}
\end{Highlighting}
\end{Shaded}

\subsection{ANEXO C.5. Ciclo de Mantenimiento y Actualización de
Versiones}\label{anexo-c.5.-ciclo-de-mantenimiento-y-actualizaciuxf3n-de-versiones}

La evolución de la librería para la corrección de vulnerabilidades o la
integración de nuevas capacidades de análisis requiere una política
estricta de control de versiones. El flujo metodológico establecido para
la liberación de actualizaciones iterativas consta de tres fases
secuenciales: - \textbf{Incremento del número de versión:} Modificación
manual del campo \texttt{version} en el archivo \texttt{pyproject.toml}
siguiendo el estándar de Versionado Semántico (ej. de \texttt{0.1.0} a
\texttt{0.1.1}). - \textbf{Saneamiento del directorio de distribución:}
Eliminación de los artefactos obsoletos del directorio \texttt{dist/}
para mitigar el riesgo de duplicidad o subidas erróneas de versiones
previas. - \textbf{Reconstrucción y despliegue:} Ejecución consecutiva
de los procesos de empaquetado y transferencia:

\begin{Shaded}
\begin{Highlighting}[]
\ExtensionTok{uv}\NormalTok{ build}
\ExtensionTok{uv}\NormalTok{ publish}
\end{Highlighting}
\end{Shaded}

Esta sistemática asegura que cada iteración de la herramienta de
auditoría de la EVM mantenga la trazabilidad, la coherencia histórica y
la disponibilidad pública necesarias para un entorno de producción
académica y profesional.
